\begin{figure}[htbp]
\centering
\begin{tikzpicture}[causal,x=1cm,y=1cm,
 ch02box/.style={draw,rounded corners=3pt,align=center,inner sep=7pt,font=\small}]
\node[ch02box,text width=7cm] (ch02root) at (0,0)
 {\textbf{机制在物理上的独立性}\\原理 2.1};
\node[ch02box,text width=3cm] (ch02left) at (-4.6,-2.7)
 {可干预性、\\自主性、\\模块性、\\不变性、\\迁移};
\node[ch02box,text width=3cm] (ch02middle) at (0,-2.7)
 {各机制所含\\信息的独立性};
\node[ch02box,text width=3cm] (ch02right) at (4.6,-2.7)
 {噪声的独立性、\\结构的\\条件独立性};
\draw[ed] (ch02root.south west)--(ch02left.north east);
\draw[ed] (ch02root.south)--(ch02middle.north);
\draw[ed] (ch02root.south east)--(ch02right.north west);
\end{tikzpicture}
\caption{机制独立原理及其对因果推断的启示（原理 2.1）。}
\label{fig:ch02-icm}
\end{figure}
